\section{Example: the integers as a ring}\label{sec:08-rings:04-integers-example:1}

We reuse \texttt{intGroup} from Chapter 6 as the additive part.

\begin{lstlisting}
def intCommGroup : CommGroup Int where
  toGroup := intGroup
  comm := by
    intro a b
    exact Int.add_comm a b
\end{lstlisting}

\texttt{toGroup := intGroup} fills the field inherited from \texttt{Group Int} inside \texttt{CommGroup Int}'s definition. This is how \texttt{extends} works mechanically: under the hood, \texttt{CommGroup G} really has fields \texttt{toGroup : Group G} and \texttt{comm : ...}, and Lean's dot-notation makes \texttt{cg.op} mean \texttt{cg.toGroup.op} automatically.

\begin{lstlisting}
def intRing : Ring Int where
  addGrp := intCommGroup
  mul := fun a b => a * b
  one := 1
  mul_assoc := by
    intro a b c
    exact Int.mul_assoc a b c
  one_mul := by
    intro a
    exact Int.one_mul a
  mul_one := by
    intro a
    exact Int.mul_one a
  left_distrib := by
    intro a b c
    exact Int.mul_add a b c
  right_distrib := by
    intro a b c
    exact Int.add_mul a b c
\end{lstlisting}

Every proof obligation is again a one-line \texttt{exact} naming a specific core-library fact about \texttt{Int} (\texttt{Int.mul\_assoc}, \texttt{Int.one\_mul}, \ldots), exactly as in Chapter 6. Integer arithmetic is not being proved from nothing; rather, already-known facts are assembled into the \texttt{Ring} bundle.

\begin{mathreading}

This shows \((\mathbb{Z}, +, \times, 0, 1)\) as an object of \(\mathbf{Ring}\) --- in fact the \hyperref[sec:01-basics:04-terminology:1]{\emph{initial} object}, since there is a unique ring homomorphism \(\mathbb{Z} \to R\) into any ring. First \texttt{intCommGroup} upgrades the additive group \((\mathbb{Z},+)\) to an abelian group by supplying commutativity (\(a + b = b + a\)); then \texttt{intRing} adds the multiplicative monoid \((\mathbb{Z}, \times, 1)\) and checks the distributive laws \(a(b+c) = ab + ac\) and \((a+b)c = ac + bc\). Each obligation is the named \(\mathbb{Z}\)-arithmetic fact, so the term is the formal counterpart of ``\(\mathbb{Z}\) is a commutative ring.''

\end{mathreading}

\textbf{Mathlib equivalent.} \texttt{Int} is already a \href{https://loogle.lean-lang.org/?q=CommRing}{\texttt{CommRing}} instance, so there is no \texttt{intRing}-style bundle to build. The obligations \texttt{intRing} checks by hand are, again, generic lemmas that hold for every commutative ring:

\begin{lstlisting}
example : CommRing Int := inferInstance

example (a b c : Int) : (a * b) * c = a * (b * c) := mul_assoc a b c
example (a : Int) : 1 * a = a := one_mul a
example (a : Int) : a * 1 = a := mul_one a
example (a b c : Int) : a * (b + c) = a * b + a * c := mul_add a b c
example (a b c : Int) : (a + b) * c = a * c + b * c := add_mul a b c
\end{lstlisting}

\href{https://loogle.lean-lang.org/?q=mul_add}{\texttt{mul\_add}}/\href{https://loogle.lean-lang.org/?q=add_mul}{\texttt{add\_mul}} are Mathlib's names for \texttt{left\_distrib}/\texttt{right\_distrib}. They are the same laws, but stated generically over \texttt{[Ring R]} (or the weaker \texttt{[Distrib R]}) instead of being cited per-type as \texttt{Int.mul\_add}/\texttt{Int.add\_mul}.
